\subsection{Proof of Main Theorem}
We work by induction on $m\ge1$.
From the definition of the non-commutative weights in equation~\eqref{eq:edge weights}, we immediately see 
\[Y_{D_1}=P_1(Y)X^{-1}=F_{P_1}(Y)=Y_1\]
and
\begin{align*}
 Y_{D_2}
 &=\sum\limits_{\ell=0}^{d_2} (P_1(Y)X^{-1})^{d_2-\ell}(X^{-1})^\ell(p_{2,d_2-\ell}X^{\ell+1}Y^{-1}X^{-1})\\
 &=P_2(P_1(Y)X^{-1})XY^{-1}X^{-1}=F_{P_1}F_{P_2}(Y)=Y_2.
\end{align*}

Write $Y_{D_2}=\sum\limits_{\omega_H:H_2\to[0,d_1]} Y_{D_2}(\omega_H)$.
We will show that $F_{P_0}(Y_{D_2}(\omega_H))=Y'_{D'_3}(\varphi^*_2\omega_H)$ for each horizontal grading $\omega_H:H_2\to[0,d_1]$.
Fix a horizontal grading $\omega_H:H_2\to[0,d_1]$.
If $\supp(\omega_H)=\emptyset$, then 
\[Y_{D_2}(\omega_H)=(X^{-1})^{d_2}P_0(X)XY^{-1}X^{-1}\]
and so 
\begin{align*}
 F_{P_0}(Y_{D_2}(\omega_H))
 &=(XY^{-1}X^{-1})^{d_2}(XP_0(Y)X^{-1})(XYX^{-1})(P_0(Y)X^{-1})^{-1}(XY^{-1}X^{-1})\\
 &=(XY^{-1}X^{-1})^{d_2}XYXY^{-1}X^{-1}\\
 &=(XY^{-1}X^{-1})^{d_2-1}X^2Y^{-1}X^{-1}\\
 &=Y'_{D'_3}(\varphi^*_2\omega_H).
\end{align*}
Suppose $\supp(\omega_H)\ne\emptyset$.
Let $h_i$ be the last horizontal edge in $\supp(\omega_H)$.
Then $(\omega_H,\omega_V)$ will be compatible if and only if $\omega_V(v_1)\le d_2-i$.
This gives
\begin{multline*}
Y_{D_2}(\omega_H)\\
=\Mk (p_{1,\omega(h_1)}Y^{\omega(h_1)}X^{-1})\Mk \cdots\Mk (p_{1,\omega(h_i)}Y^{\omega(h_i)}X^{-1})(X^{-1})^{d_2-i}\sum\limits_{\ell=0}^{d_2-i} p_{2,d_2-\ell}X^{\ell+1}Y^{-1}X^{-1}.
\end{multline*}
Applying $F_{P_0}$ gives
\begin{multline*}
 F_{P_0}(Y_{D_2}(\omega_H))\\
\begin{aligned}
 &=\big(p_{1,\omega(h_1)}(P_0(Y)X^{-1})^{\omega(h_1)}XY^{-1}X^{-1}\big)\cdots\big(p_{1,\omega(h_i)}(P_0(Y)X^{-1})^{\omega(h_i)}XY^{-1}X^{-1}\big)\\
 &\quad\times[XY^{-1}X^{-1}]^{d_2-i}\sum\limits_{\ell=0}^{d_2-i} p_{2,d_2-\ell}(XYX^{-1})^{\ell+1}(P_0(Y)X^{-1})^{-1}(XY^{-1}X^{-1})\\
 &=\Mk(\mk P_0(\mk Y\mk )X^{\mk -\mk 1}\mk )^{\omega(\mk h_1\mk ) }\Mk (\mk X^{\mk -\mk 1}\mk )^{d_1\mk -\mk \omega(\mk h_1\mk )}(\mk p_{1,\omega(h_1\mk )}X^{d_1\mk -\mk \omega(\mk h_1\mk )\mk +\mk 1}Y^{\mk -\mk 1}X^{\mk -\mk 1})\Mk \cdots\Mk (\mk P_0(\mk Y\mk )X^{\mk -\mk 1}\mk )^{\omega(\mk h_i\mk )\mk -\mk 1}\\
 &\quad\times\left(\sum\limits_{\ell=0}^{d_2-i} p_{2,d_2-\ell} Y^\ell X^{-1}\right)(X^{-1})^{d_1-\omega(h_i)}(p_{1,\omega(h_i)}X^{d_1-\omega(h_i)+1}Y^{-1}X^{-1})\\
 &\quad\big[(X^{-1})^{d_1}(X^{d_1+1}Y^{-1}X^{-1})\big]^{d_2-i-1}(X^{-1})^{d_1-1}(X^{d_1+1}Y^{-1}X^{-1})\\
 &=Y'_{D'_3}(\varphi^*_2\omega_H).
\end{aligned}
\end{multline*}

Suppose $m\ge3$ and let $\omega_H:H_m\to[0,d_1]$ be a horizontal grading of $D_m$.
Following Theorem~\ref{th:blocking edge conditions}, there are two cases to consider.
\begin{enumerate}[label=(\alph*)]
 \item Suppose that $(\omega_H,\omega_V)$ is compatible for every piecewise compatible grading $(\omega_H,\omega_V)$ of $D_m$.
 Then Lemma~\ref{le:piecewise compatible factorization 1} shows there is the factorization
 \begin{multline}\label{dagger3}
 \tag{$\dagger$}
 Y_{D_m}(\omega_H)\\
 =\Mk Y_{D_{m\mk -\mk 1}}(\omega_{H,1})Y_{D_{m\mk -\mk 1}}(\omega_{H,2})\Mk\cdots\Mk Y_{D_{m\mk -\mk 1}}(\omega_{H,d_m\mk -\mk 1\mk -\mk \delta_m})pX^{|H_{m\mk -\mk 2\mk -\mk \delta_m}|} Y_{D_{m-1}}(\omega_{H,d_m-\delta_m}),
 \end{multline}
 where 
 \[p=\begin{cases} p_{1,1}^{|V_{m-2-\delta_m}|-2|H_{m-2-\delta_m}|}p_{1,2}^{|H_{m-2-\delta_m}|-|V_{m-2-\delta_m}|} & \text{if $d_2=1$;}\\ p_{1,1}^{-|V_{m-2-\delta_m}|} & \text{if $d_2>1$.}\end{cases}\]
 If $m\ge3+\delta'_{m+1}$, we may apply Lemma~\ref{le:piecewise compatible factorization 2} to conclude by induction that $F_{P_0}(Y_{D_m}(\omega_H))=Y'_{D'_{m+1}}(\varphi^*_m\omega_H)$.

 It remains to consider the case $m=3$ with $\delta'_4=1$, \ie $d_2=1$.
 In this case $D_2$ consists of a single horizontal edge followed by a single vertical edge and, by Corollary~\ref{cor:recursive dyck path structure}\ref{coro2.4_3}, $D_3$ consists of $d_1-1$ copies of $D_2$ followed by a single vertical edge, in particular $D_3$ ends with two consecutive vertical edges.
 The factorization~\eqref{dagger3} still holds in this case and by induction we have $F_{P_0}(Y_{D_2}(\omega_{H,r}))=Y'_{D'_3}(\varphi^*_2\omega_{H,r})$ for $1\le r\le d_1$.
 In particular, by Lemma~\ref{le:piecewise compatible factorization 3} to see that $F_{P_0}(Y_{D_3}(\omega_H))=Y'_{D'_4}(\varphi^*_3\omega_H)$ it suffices to compare $F_{P_0}\big(Y_{D_2}(\omega_{H,d_1-1})pXY_{D_2}(\omega_{H,d_1})\big)$ with $XY'_{D'_3}(\omega_{V',d_1-1})$, where we write $\omega_{V'}=\varphi^*_3\omega_H$.

 There are two cases to consider.
 If $\omega(h_{d_1-1})=0$, we have $Y_{D_2}(\omega_{H,d_1-1})=(X^{-1})P_0(X)XY^{-1}X^{-1}$ and so
 \begin{align*}
 F_{P_0}(Y_{D_2}(\omega_{H,d_1-1}))
 &\Mk =\Mk (XY^{-1}X^{-1})(XP_0(Y)X^{-1})(XYX^{-1})(P_0(Y)X^{-1})^{-1}XY^{-1}X^{-1}\\
 &\Mk =\Mk X^2Y^{-1}X^{-1}.
 \end{align*}
 The same calculation shows $F_{P_0}\big(Y_{D_2}(\omega_{H,d_1})\big)=X^2Y^{-1}X^{-1}$ by the assumptions on $\omega_{H,d_1}$ in Definition~\ref{def:agregate grading}.
 But then, since $p=1$ in this case, we have
 \begin{align*}
 F_{P_0}\big(Y_{D_2}(\omega_{H,d_1-1})pXY_{D_2}(\omega_{H,d_1})\big)
 &=(X^2Y^{-1}X^{-1})(XYX^{-1})(X^2Y^{-1}X^{-1})\\
 &=X^3Y^{-1}X^{-1}\\
 &=X(X^{-1})^{d_1-1}(X^{d_1+1}Y^{-1}X^{-1}),
 \end{align*}
 which is exactly $XY'_{D'_3}(\omega_{V',d_1-1})$.
 
 When $h_{d_1\mk -\mk 1}\Mk \in\Mk \supp(\omega_{\mk H})$, we have $Y_{D_2}(\omega_{H,d_1\mk -\mk 1})\Mk =\Mk p_{1,\omega(h_{d_1\mk -\mk 1}}Y^{\omega(h_{d_1\mk -\mk 1})\mk -\mk 1}X^{\mk -\mk 1}$ so that
% \[
\begin{multline*}
 F_{P_0}\big(Y_{D_2}(\omega_{H,d_1-1})\big)
 \\
 =(P_0(Y)X^{-1})^{\omega(h_{d_1-1})-1}(X^{-1})^{d_1-\omega(h_{d_1-1})}p_{1,\omega(h_{d_1-1})}X^{d_1-\omega(h_{d_1-1})+1}Y^{-1}X^{-1}.
 %\]
 \end{multline*}
 We saw above that $F_{P_0}\big(Y_{D_2}(\omega_{H,d_1})\big)=X^2Y^{-1}X^{-1}$ and so
 \begin{align*}
 &F_{P_0}\big(Y_{D_2}(\omega_{H,d_1-1})pXY_{D_2}(\omega_{H,d_1})\big)\\
 &\qquad=(P_0(Y)X^{-1})^{\omega(h_{d_1-1})-1}(X^{-1})^{d_1-\omega(h_{d_1-1})-1}p_{1,\omega(h_{d_1-1})}X^{d_1-\omega(h_{d_1-1})+1}Y^{-1}X^{-1},
 \end{align*}
 which is exactly $XY'_{D'_3}(\omega_{V',d_1-1})$.
 The claim then follows by induction from Lemma~\ref{le:piecewise compatible factorization 3}.
 \item Suppose there exists a vertical grading $\omega_V^*:V_m\to[0,d_2]$ of $D_m$ so that $(\omega_H,\omega_V^*)$ is piecewise compatible, but not compatible. 
 By Theorem~\ref{th:blocking edge conditions}, there must exist a blocking edge $h_i$ for $\omega_H$.
 Set 
 \[d=|\supp(\omega_H)\cap h_iv_{u_{m-1,2}}|\quad\text{ and }\quad t=|\supp(\omega_V^*)\cap h_iv_{u_{m-1,2}}|.\]

 By Proposition~\ref{prop:compatibility equivalence} and Proposition~\ref{prop:piecewise equivalence}, $(\Omega_m\omega^*_V,\varphi^*_m\omega_H)$ is a piecewise compatible grading of $D'_{m+1}$ which is not compatible.
 Let $D(v'_{u'_{m,2}};\varphi^*_m\omega_H)=h'_jv'_{u'_{m,2}}$ and observe that $\hgt(h'_j)=i-1$ by definition of $\Omega_m$.
 By Proposition~\ref{prop:strong justification implication}, Lemma~\ref{le:remote shadow cardinalities}, and Corollary~\ref{cor:support and remote shadow}, we have
 \begin{align*}
 |\supp(\Omega_m\omega^*_V)\cap h'_jv'_{u'_{m,2}}|&=|\rsh(\varphi^*_m\omega_H)\cap h'_jv'_{u'_{m,2}}|=|\rsh(\omega_H)\cap h_iv_{u_{m-1,2}}|\\
 &=|\supp(\omega_V)\cap h_iv_{u_{m-1,2}}|.
 \end{align*}
 Moreover, we have
 \[|\supp(\varphi^*_m\omega_H)\cap h'_jv'_{u'_{m,2}}|=u_{m,1}-i+1-|\supp(\omega_H)\cap h_iv_{u_{m-1,2}}|+\delta,\]
 where $\delta=0$ if $\omega_H(h_{i-1+d})=d_1$ and $\delta=1$ otherwise.
 If then follows from the definitions of $\varphi^*_m$ and $\Omega_m$ that the coefficients $p_2$ agree in Lemma~\ref{le:incompatible factorization 1} and Lemma~\ref{le:incompatible factorization 2}.

 Using the notation of Lemma~\ref{le:piecewise compatible factorization 1} and Lemma~\ref{le:incompatible factorization 1}, we have 
 \[
 Y_{D_m}(\omega_H)=Y_{D_m}^{pc}(\omega_H)-Y_{D_m}^{nc}(\omega_H).
 \]
 By induction we have $F_{P_0}(Y_{D_{m-1}}(\omega_{H,r}))=Y'_{D'_m}(\varphi^*_{m-1}\omega_{H,r})$ for $1\le r\le d_m-\delta_m$ and $F_{P_0}(Y_{D_{m-1}}(\chi_H))=Y'_{D'_m}(\chi_{V'})$.
 It follows that 
 \[
 F_{P_0}\big(Y_{D_m}^{pc}(\omega_H)\big)={Y'}_{D'_{m+1}}^{pc}(\varphi^*_m\omega_H) 
 \quad\text{and} \quad F_{P_0}\big(Y_{D_m}^{nc}(\omega_H)\big)={Y'}_{D'_{m+1}}^{nc}(\varphi^*_m\omega_H).
 \]
 Since $Y'_{D'_{m+1}}(\varphi^*_m\omega_H)={Y'}_{D'_{m+1}}^{pc}(\varphi^*_m\omega_H)-{Y'}_{D'_{m+1}}^{nc}(\varphi^*_m\omega_H)$, the result follows.
\end{enumerate}
